\documentclass{article}
\usepackage[T2A]{fontenc}
\usepackage[utf8]{inputenc} %это кодировка UTF-8 точно работает
\usepackage{amsmath}
\usepackage{amssymb}
\usepackage{amsfonts}
\usepackage{mathrsfs}
\usepackage[12pt]{extsizes}
\usepackage{fancyvrb}
\usepackage{indentfirst}
\usepackage[
left=2cm, right=2cm, top=2cm, bottom=2cm, headsep=0.2cm, footskip=0.6cm, bindingoffset=0cm
]{geometry}
\usepackage[english,russian]{babel}

\begin{document}

\section*{Вариант 26}

Уравнение математической модели:

 \begin{equation}
    q_0 \left( q_0^0 + q_T^0 + K p_x^0 \right) + 
    q \left( q_0^T + q_T^T + K p_x^T \right) + 
    p_0 \left( q_0^{TP} + q_T^{TP} + K p_x^{TP} \right) = 1.
 \end{equation}

По аналогии с предыдущим вариантом модели условные вероятности 
$q_0((l+1)/l)$, $q((l+1)/l)$, $p_x((l+1)/l)$ правильного приема нулевого, 
токового, стирания, соответственно, для $(l+1)$-го символа запишутся в виде:

\begin{equation}
    \left.
    \begin{aligned}
        q_0((l+1)/l) &= q_0 q_0^0 + q q_0^T + K p_x q_0^x; \\
        q_T((l+1)/l) &= q_0 q_T^0 + q q_T^T + K p_x q_T^x; \\
        p_0((l+1)/l) &= q_0 p_TP^0 + q p_{TP}^T + K p_x p_x^x.
    \end{aligned}
    \right\}
\end{equation}

Решая систему уравнений (2), можно получить формулы для определения безусловных вероятностей $q_0$, $q$, $p_x$.

\end{document}
